\subsection{Inconsistent and Dependent Systems} 
The solution to a system of two equations in two variables represents the point of intersection of two lines. What can go wrong?
\begin{example}
Consider the system:
\begin{syspic}{3}
\draw (-.5\textwidth,0) node{%
$
\begin{system}
3x+2y=6\\
y=-\frac{3}{2}x-1
\end{system}
$};
\draw[graph] (0,3) -- (3,-3/2);
\draw[graph] (-8/3,3) -- (4/3,-3);
\end{syspic}
These lines are \makeblank[1in] so they \makeblank. This system has \makeblank. If we try to solve it:
\begin{lpic}{}{5}
\twocolleft{-1}{5}{Substitution:}{Elimination:}
\end{lpic}
These equations contain no \makeblank[1in] and are \makeblank[1in]. A system like this is called \makeblank. Notice that in standard form, these are obviously \makeblank[1.5in] lines.
\end{example}
%%%%%%%%%%%%%%%%%%%%%%%%%%%%%%%%%%%%%%%%%%%%%%%%%%%%%%%%%%%
\begin{example}
Consider the system:
\begin{syspic}{3}
\draw (-.5\textwidth,0) node{%
$
\begin{system}
3x-4y=4\\
y=\frac{3}{4}x-1
\end{system}
$};
\draw[graph] (-8/3,-3) -- (3,5/4);
\end{syspic}
These two equations' graphs are \makeblank. This system has \makeblank. If we try to solve it:
\begin{lpic}{}{5}
\twocolleft{-1}{5}{Substitution:}{Elimination:}
\end{lpic}
These equations contain no \makeblank[1in] and are \makeblank[1in]. A system like this is called \makeblank. Notice that in standard form, these are obviously \makeblank[1.5in] line.
\end{example}
